\section[title={2024-09-18 - Vergleichen der beiden
Lebenskonzepte},reference={2024-09-18-Vergleichen der beiden
Lebenskonzepte}]

\startframedtask
{\bf Vergleichen} Sie die beiden gegensätzlichen Lebensentwürfen
mithilfe geeigneter Begriffspaare wie etwa Ordnung / Unornung, Stadt /
Land, Rationalität / Emotionalität usw. Nutzen Sie hierfür insbesondere
die Seiten 11f., 19f., 22f., 34f., 38f., 60f., 68f., 71f., 91f., 103f.,
111f., 120f.

Meine Gruppe ab 68f.

\stopframedtask

\startitemize[packed]
\item
  S. 71
  \startitemize[packed]
  \item
    Empfindlichkeit gegen die Sonne, ist abhängig von der
    Characterstärke
  \stopitemize
\stopitemize

\page

\startcolumnset[n=2]

\startsectionlevel[title={Eine \quotation{richtige}
Familie},reference={eine-richtige-familie}]

\startsectionlevel[title={Ordnung},reference={ordnung}]

\startitemize[packed]
\item
  Stille
\item
  Perfektionismus
\item
  Tagesstruktur
\item
  Unterordnung
\item
  Psychischer Druck gegenüber Familie
\stopitemize

\stopsectionlevel

\startsectionlevel[title={Stadt},reference={stadt}]

\startitemize[packed]
\item
  Sozialer Aufstieg / Beförderung
\item
  Symbol für Ordnung / Struktur / Fortschritt
\item
  Wohlstand
\stopitemize

\stopsectionlevel

\startsectionlevel[title={Rationalität},reference={rationalität}]

\startitemize[packed]
\item
  gegen Emotionalität / bekämpft sie mit Rationalität
\item
  Rationalität ist der Schlüssel zum Erfolg
\item
  Vater ist das Optimum; Unterbindet die
\stopitemize

\column

\stopsectionlevel

\stopsectionlevel

\startsectionlevel[title={\quotation{Verwildert}},reference={verwildert}]

\startsectionlevel[title={Unordnung},reference={unordnung}]

\startitemize[packed]
\item
  Gleichgültigkeit / Gelassenheit als Wunsch für familiären frieden
\item
  Freiheit = Ordnung (Individuale Entfaltung als Ziel)
\stopitemize

\stopsectionlevel

\startsectionlevel[title={Land},reference={land}]

\startitemize[packed]
\item
  Stadt ist \quotation{großer Käfig}
\item
  {\em Land} als Symbol für Freiheit und Persöhnlichkeitsentwicklung
\item
  Symbol der Still
\item
  Möglichkeit der freien Bildung
\item
  Freie Kunst
\item
  Ort der schönen Geister
\stopitemize

\stopsectionlevel

\startsectionlevel[title={Emotionalität},reference={emotionalität}]

\startitemize[packed]
\item
  Emotionalität wird mit Leichtsinnigkeit gleichgesetzt
  \startitemize[packed]
  \item
    wurde von der Tocher, von dem Vater, übernommen
  \stopitemize
\item
  Emotionalität gefährdet die Vormachtstellung des Vaters
\stopitemize

\stopcolumnset

\stopsectionlevel

\stopsectionlevel

\startsectionlevel[title={Hausaufgabe: Umgang mit Konflikten in der
Familie},reference={hausaufgabe-umgang-mit-konflikten-in-der-familie}]

\startsectionlevel[title={Mutter},reference={mutter}]

\startitemize[packed]
\item
  Stellt sich um
\stopitemize

\stopsectionlevel

\startsectionlevel[title={Tochter},reference={tochter}]

\startitemize[packed]
\item
  Rebelling / Lehnt sich auf
\stopitemize

\stopsectionlevel

\startsectionlevel[title={Sohn:},reference={sohn}]

\startitemize[packed]
\item
  Singen
\item
  Suizidal
\stopitemize

\stopsectionlevel

\stopsectionlevel
